% =============================================================================
% D8 — Financement et cycles d'infrastructure — v2 (mai 2026)
% Réécriture après diagnostic : objet = cycle macrofinancier de l'investissement
% en infrastructure informationnelle. Pas de matériau D7 (transaction costs).
% =============================================================================


\subsection{Financement et cycles d'infrastructure}\label{sec:d8-finance}

\subsubsection{Le champ~: le cycle d'investissement en infrastructure informationnelle}

Le chapitre~1 a documenté une régularité que la théorie économique standard ne prédit pas~: la capacité technique de l'infrastructure informationnelle précède systématiquement son déploiement, parfois de plusieurs décennies. Le câble transatlantique était techniquement réalisable dès~1857~; il fallut cinq tentatives, deux faillites et neuf ans pour qu'un câble fonctionnel soit posé \parencite{gordon_thread_2002}. La machine analytique de Babbage était constructible dans les années~1840 \parencite{swade_difference_2001}~; elle ne fut jamais achevée. Le réseau électrique d'Insull produisait un kilowattheure moins cher que les dynamos privées dès les années~1890~; il ne se généralisa qu'avec l'invention d'un cadre réglementaire qui garantissait à l'investisseur un rendement en échange d'une obligation de service \parencite{hughes_networks_1983}. L'\emph{Electric Lighting Act} de~1882, à Londres, gela l'investissement privé pendant vingt ans par une clause de rachat municipal à la valeur de la ferraille (\emph{scrap iron clause})~: soixante-neuf autorisations en~1883, zéro en~1885. Ce n'est pas la technique qui fait défaut. C'est le véhicule financier.

\markNEO Le théorème de \textcite{modigliani_cost_1958} pose le benchmark~: sous des conditions idéales, la structure de financement n'affecte pas l'allocation du capital. \NEO{Si Modigliani-Miller tenait, la dimension~8 n'existerait pas, et le déploiement ne dépendrait que de la technologie et de la demande.} Or l'infrastructure informationnelle viole chacune des hypothèses du théorème~--- coûts fixes irréversibles, revenus incertains et différés, asymétrie informationnelle entre investisseur et opérateur, marchés incomplets pour le risque technologique. Ces violations ne sont pas des imperfections marginales~; elles sont constitutives du type d'actif. Dès lors, la forme du financement --- concession, location avec consommable, régulation des \emph{utilities}, venture capital, project finance --- n'est pas un détail d'exécution~; c'est un déterminant de ce qui se construit et de ce qui ne se construit pas.

La question analytique de D8 se formule en deux volets~: la séquence investissement massif, crise, recapitalisation est-elle récurrente dans l'infrastructure informationnelle~? Et le véhicule financier se transforme-t-il avec le type de capital~?


\subsubsection{Trois mécanismes causaux}

\paragraph{Premier mécanisme --- Le cycle investissement-crise-transfert}

\markEVO \textcite{perez_technological_2002} a proposé le cadre le plus abouti pour ce mécanisme. \EVO{Chaque révolution technologique traverse quatre phases --- irruption, frénésie, synergie, maturité --- et la transition de la frénésie à la synergie passe par une crise financière qui n'est pas un accident mais un mécanisme de transfert~: le capital spéculatif finance l'installation de l'infrastructure~; la crise en transfère la propriété à des opérateurs qui la déploient à un coût inférieur au coût de construction.} \textcite{reinhart_this_2009} confirment le pattern par un autre registre~: le syndrome \og \emph{this time is different}\fg{}, documenté sur huit siècles de crises financières, accompagne systématiquement les phases d'installation.

Le mécanisme a trois temps. Dans le premier, les rendements croissants du réseau (D2) attirent un capital dont le volume dépasse ce que le calcul prudent justifierait~: investir est rationnel si l'on croit être dans le camp du vainqueur, et les rendements croissants font du marché un jeu à un seul gagnant. Dans le deuxième, la logique \emph{hedge} au sens de \textcite{minsky_stabilizing_1986} --- où les revenus couvrent le service de la dette --- glisse vers une logique spéculative, puis éventuellement Ponzi~: la continuation du financement dépend de la hausse des valorisations, non des flux de trésorerie. Dans le troisième, la crise transfère l'infrastructure à des opérateurs qui l'acquièrent à prix réduit et la rendent profitable.

Le capital spéculatif joue un rôle fonctionnel que Perez identifie avec précision~: il prend des risques que le capital productif ne prendrait pas, parce que son horizon est plus court et sa capacité à sortir plus grande. La bulle des canaux (1790s), la \emph{railway mania} (1840s), la \emph{corporate mania} électrique (1880s--1890s) et la bulle dotcom (2000) instancient toutes la même séquence. Le capital détruit dans la crise n'est pas perdu pour l'économie~; l'infrastructure survit à ses financeurs. La fibre optique posée pendant la bulle dotcom a servi de substrat à l'économie numérique de la décennie suivante.

Les conditions d'activation sont~: rendements croissants de réseau, coûts fixes élevés et irréversibles, asymétrie informationnelle entre investisseurs et opérateurs, et liquidité financière suffisante pour alimenter la frénésie. Quand l'une manque, le mécanisme ne s'enclenche pas. Le marché des calculatrices mécaniques, sans rendements de réseau ni coûts fixes prohibitifs, n'a connu ni bulle ni consolidation~: une centaine de firmes coexistaient en~1904 \parencite{cortada_before_1993}.


\paragraph{Deuxième mécanisme --- L'adéquation du véhicule à la structure temporelle}

Le premier mécanisme porte sur le volume et le cycle du capital~; le deuxième porte sur sa forme. Le déploiement échoue quand le véhicule de financement est inadapté à la structure temporelle des coûts et des revenus de l'infrastructure, même quand la technique fonctionne et la demande existe.

Toute infrastructure crée un décalage entre le moment où les coûts sont engagés --- ponctuel et irréversible --- et le moment où les revenus se matérialisent --- progressif et incertain. Le véhicule financier résout ce décalage~: il convertit une dépense immédiate en un flux de revenus futur que l'investisseur accepte de financer. La forme du véhicule dépend de la nature de l'actif.

\markINFO \textcite{hall_financing_2010} formalisent les contraintes~: incertitude sur les résultats, intangibilité de l'actif, asymétrie informationnelle. \INFO{Ces propriétés créent un \og \emph{financing gap}\fg{} qui biaise l'allocation~: les projets à rendements sociaux élevés mais privément non appropriables sont sous-financés.} Mais le \emph{financing gap} de Hall et Lerner est uniforme~; ce que le matériau du chapitre~1 montre, c'est que la contrainte varie avec le type de capital (D3). Un câble sous-marin est tangible, collatéralisable et à revenus prévisibles sur un corridor dense~; un algorithme est intangible et non collatéralisable~; une tabulatrice est un bien d'équipement coûteux mais dont les cartes perforées fournissent un flux récurrent. Chaque configuration appelle un véhicule distinct~:

La \textbf{concession} --- où l'État accorde un monopole temporaire contre l'obligation de construire --- est le véhicule du câble transatlantique et des premiers réseaux télégraphiques. Elle fonctionne quand la demande est avérée et concentrée sur un corridor~; elle échoue quand la demande est hypothétique (câbles transpacifiques de~1902, \textcite{muller_telegraph_2015}). La \textbf{location avec consommable exclusif} est le véhicule d'IBM~: la tabulatrice à 40~dollars par mois et les cartes à 85~cents le millier transforment un \emph{capex} en \emph{opex} récurrent à marge élevée (60 à 80\,\%), et IBM paya des dividendes sans interruption pendant la Grande Dépression \parencite{cortada_before_1993}. La \textbf{régulation des \emph{utilities}} est le véhicule d'Insull~: le régulateur garantit un rendement sur capital investi en échange d'une obligation de service, transformant l'actif irréversible en quasi-bond~; les \emph{public utility commissions}, créées à partir de~1907 aux États-Unis, institutionnalisent ce contrat \parencite{hughes_networks_1983}. Le \textbf{brevet de principe} est le véhicule du téléphone de Bell~: le juge Lowell (1881) qualifie la transmission de la parole de \og \emph{new art}\fg{} et accorde la revendication la plus large, produisant un rendement de quarante-six pour cent pendant dix-huit ans \parencite{gabel_early_1969}. Le \textbf{venture capital} est le véhicule de l'intangible scalable~: il résout l'asymétrie informationnelle par le contrôle actif, la sélection et le financement par étapes \parencite{hall_financing_2010}.

Ce qui donne au mécanisme sa force prédictive, c'est le test en négatif. Quand le véhicule est inadapté, la technique échoue commercialement. Babbage disposait d'un financement public unique (17\,470~livres) pour un projet dont les coûts étaient cumulatifs et les revenus inexistants~: pas de flux, pas de recapitalisation, pas de déploiement. Le câble de~1857 fit faillite en trois semaines~; mais la promesse d'arbitrage intercontinental --- un flux identifiable --- permit la recapitalisation de~1864. À Londres, la \emph{scrap iron clause} de~1882 offrait aux investisseurs un rachat municipal à la valeur des composants séparés, sans prime pour la valeur du système en fonctionnement~: soixante-neuf autorisations provisoires en~1883, zéro en~1885, la station Edison de Holborn Viaduct abandonnée en~1886 \parencite{hughes_networks_1983}. Le véhicule ne détruisait pas l'infrastructure~; il détruisait l'incitation à la construire. Le capital ne finance pas une technique~; il finance un flux de revenus. Pas de flux, pas de déploiement.


\paragraph{Troisième mécanisme --- L'État comme preneur de risque de premier rang}

Quand le risque est trop élevé ou trop incertain pour que le capital privé l'absorbe, l'État finance la première génération de l'infrastructure, socialise le risque de développement, et le secteur privé capte les rendements dans la phase de déploiement. \textcite{mazzucato_entrepreneurial_2013} a systématisé ce mécanisme pour la période contemporaine~; le matériau du chapitre~1 montre qu'il est structurel dans l'histoire de l'infrastructure informationnelle.

En~1838, le Congrès américain finança la première ligne télégraphique Washington-Baltimore après six ans de débat parlementaire~: la communication instantanée relevait-elle de l'utilité générale~? La ligne fonctionna~; le réseau se développa ensuite en concurrence privée anarchique, puis en monopole Western Union \parencite{thompson_wiring_1947}. Le gouvernement britannique finança Babbage à hauteur de 17\,470~livres entre~1823 et~1842~; le projet échoua, mais le risque fut intégralement porté par le contribuable. À Bletchley Park, dix mille personnes travaillèrent au déchiffrement d'Enigma avec des ressources sans limites ordinaires~; Colossus (1943) fut le premier ordinateur électronique programmable opérationnel, mais le secret militaire empêcha le transfert direct. Les ingénieurs de Bletchley furent recrutés par IBM, Ferranti et GEC~; les savoirs dont le coût avait été supporté par le contribuable alimentèrent l'industrie informatique d'après-guerre.

Le mécanisme a une structure temporelle caractéristique. L'État intervient dans la phase où le rapport risque/rendement est inacceptable pour le capital privé~: soit parce que les revenus sont trop incertains (le télégraphe en~1838), soit parce que la demande est militaire et non monétisable (Bletchley Park), soit parce que l'horizon dépasse celui du capital privé (les infrastructures de recherche fondamentale). Une fois la preuve de concept établie et le risque réduit, le capital privé prend le relais~--- souvent en acquérant à bas prix des actifs dont le développement a été financé publiquement.

\textcolor{teal}{Ce mécanisme a une implication directe pour le cycle contemporain. Le CHIPS and Science Act (2022) aux États-Unis et le European Chips Act (2023) en Europe sont des subventions publiques massives à la fabrication de semi-conducteurs. L'argument de Mazzucato prédit que les rendements de ces investissements seront captés par les firmes privées qui en bénéficient. La question est de savoir si les conditions de ces subventions intègrent un mécanisme de retour vers le financeur public, ou si le pattern historique --- socialisation du risque, privatisation du rendement --- se reproduit. La question est posée~; elle n'est pas tranchée.}


\subsubsection{Confrontation à l'évidence}

\paragraph{Évidence historique}

Les trois mécanismes s'instancient selon des combinaisons distinctes dans les cas du chapitre~1, ce qui permet de vérifier les conditions d'activation.

\emph{Premier temps~: le cycle spéculatif (M1).} Le télégraphe américain suit le pattern Perez complet. Vingt compagnies en concurrence anarchique dans les années~1850, consolidation sous Western Union en~1866, monopole profitable dans les années~1870 \parencite{nonnenmacher_telegraph_2001}. Le câble transatlantique suit la même séquence en accéléré~: le premier câble de~1857--58 coûte 300\,000~livres à l'\emph{Atlantic Telegraph Company}, ne fonctionne que trois semaines, et la compagnie fait faillite~; la recapitalisation de~1864 finance un câble réussi en~1866 pour 600\,000 à 700\,000~livres \parencite{gordon_thread_2002}. Le téléphone confirme la séquence par un troisième canal~: le monopole de brevet (1876--1894) produit quarante-six pour cent de rendement et la stagnation du réseau~; l'expiration du brevet ouvre une concurrence explosive (six mille compagnies en~1900)~; la crise des profits conduit le syndicat Morgan-Vail à racheter et reconsolider le réseau à partir de~1907 \parencite{mueller_universal_1993}. Dans chaque cas, l'infrastructure survit à la crise~; ce sont les financeurs initiaux qui sont éliminés.

\emph{Deuxième temps~: le véhicule adapté et l'anti-pattern (M2).} Le contraste Babbage/câble est le test le plus net. Babbage avait une technique fonctionnelle et un financement public unique~; le câble avait une technique identique en~1857 et en~1866. Ce qui sépare l'échec du succès, c'est la présence d'un flux de revenus identifiable (l'arbitrage intercontinental) qui attire du capital neuf après la faillite. Babbage n'avait que la promesse scientifique~; le gouvernement britannique, après dix-neuf ans, ne la jugea plus suffisante \parencite{swade_difference_2001}. IBM fournit le test positif symétrique~: le modèle location-consommable transforme un bien d'équipement coûteux en un flux récurrent à marge élevée, et l'entreprise traverse la Grande Dépression sans interruption de dividendes \parencite{cortada_before_1993}. À l'inverse, la \emph{scrap iron clause} londonienne montre qu'un véhicule mal conçu peut geler l'investissement dans une infrastructure dont la demande existe~: en~1918, Londres comptait encore soixante-cinq compagnies d'électricité, dix fréquences différentes et vingt-quatre tensions, alors que Berlin, dont la concession municipale garantissait la stabilité de l'investissement, avait unifié son réseau avant~1914 \parencite{hughes_networks_1983}.

\emph{Troisième temps~: l'État financeur (M3).} Le Congrès américain finance la première ligne télégraphique en~1844 après six ans d'obstruction parlementaire. Le gouvernement britannique finance Babbage entre~1823 et~1842. L'État français nationalise le réseau téléphonique en~1889 et le finance par les \og avances remboursables\fg{} des collectivités locales --- un véhicule qui fragmente le réseau selon la géographie administrative plutôt que selon la rationalité économique \parencite{bertho_histoire_1984}. Bletchley Park produit Colossus avec des ressources militaires illimitées~; le secret empêche le transfert, mais les ingénieurs migrent vers l'industrie privée. Dans chaque cas, l'État absorbe le risque que le capital privé ne peut pas évaluer~; dans chaque cas, la phase de déploiement est privée.


\paragraph{Évidence contemporaine}

Le marqueur financier le plus net de la période contemporaine est le passage du venture capital au project finance comme véhicule dominant de l'investissement en infrastructure informationnelle. Pendant trente ans (1990--2020), le secteur numérique a été financé par le VC et par l'autofinancement des grandes firmes technologiques. L'émergence des data centers à l'échelle du gigawatt et des contrats d'achat d'énergie de long terme marque un basculement~: l'infrastructure informationnelle redevient tangible, et avec elle le véhicule financier bascule.

Ce basculement instancie le deuxième mécanisme. L'entraînement d'un modèle de langage coûte plusieurs centaines de millions de dollars en \emph{compute} et nécessite un data center dont la durée de vie utile se mesure en années. Le capital-risque, dont l'horizon est de cinq à sept ans et dont la logique est celle de la sortie par cession, est structurellement inadapté au financement d'un actif dont le coût de fonctionnement est permanent. Le project finance est la réponse~: le contrat de long terme garantit un flux de revenus qui sert de collatéral, et le risque est porté par le bilan de l'acheteur plutôt que par le marché des actions.

\textcolor{red}{L'analogie avec le schéma de Perez appliquée au cycle d'investissement IA est tentante~: la phase actuelle (2023--2026) présente les symptômes de la frénésie. Mais l'objection est sérieuse. Les firmes qui financent l'infrastructure disposent de cash-flows opérationnels considérables et financent une grande partie de l'investissement sur fonds propres, pas par endettement spéculatif. Le risque n'est pas la faillite~; c'est l'allocation massive de capital dans une infrastructure dont le coût de fonctionnement est permanent et dont les revenus ne sont pas garantis. La question de savoir si le cycle actuel suivra le pattern Perez complet ou s'il représente un cas nouveau --- un investissement autofinancé dont le risque est l'allocation sous-optimale plutôt que la faillite --- reste ouverte.}

\textcolor{teal}{Le troisième mécanisme est actif dans la période contemporaine. Le CHIPS and Science Act (2022, 52~milliards de dollars) et le European Chips Act (2023) reproduisent le pattern de Bletchley Park et du Congrès de~1838~: l'État finance l'infrastructure dont le risque excède la capacité du capital privé. Les subventions publiques aux usines de semi-conducteurs sont le véhicule financier du moment~; la question de Mazzucato --- le retour vers le financeur public est-il intégré dans les conditions --- se pose avec la même acuité qu'en~1944.}

\wip{Données manquantes~: (1)~les données de Miller sur l'investissement au pic ($\sim$6\,\% du PIB, identique railways/ICT/cycle en cours) doivent être sourcées précisément~; (2)~les montants de capex des hyperscalers (Microsoft, Google, Amazon, Meta) sur la période 2023--2026 et leur structure de financement (autofinancement vs.\@ dette) sont nécessaires pour fonder empiriquement le basculement VC $\rightarrow$ project finance~; (3)~les rapports IEA et OECD sur les investissements en data centers et infrastructure énergétique fourniraient le matériau quantitatif qui manque à cette section.}


\subsubsection{Fait stylisé}

L'infrastructure informationnelle ne se déploie pas sans un véhicule de financement adapté à sa structure de coûts. Le capital ne finance pas une technique~; il finance un flux de revenus. Quand ce flux est identifiable et appropriable (l'arbitrage intercontinental du câble, les cartes perforées d'IBM, le kilowattheure régulé d'Insull), le déploiement procède. Quand il est absent (Babbage) ou quand le véhicule le détruit (la \emph{scrap iron clause} londonienne), la technique échoue commercialement. Le déploiement passe par un cycle récurrent --- investissement spéculatif, crise, transfert d'infrastructure --- documenté pour le télégraphe, les chemins de fer, l'électricité et les dotcoms \parencite{perez_technological_2002, nonnenmacher_telegraph_2001, reinhart_this_2009}. Le véhicule financier se transforme avec le type de capital~: concession pour les réseaux tangibles, location-consommable pour les systèmes intégrés, régulation des \emph{utilities} pour les monopoles naturels, venture capital pour les actifs intangibles scalables, financement public pour le risque que le capital privé ne sait pas évaluer. Le passage contemporain du venture capital au project finance est le marqueur financier du retour de la tangibilité dans l'infrastructure informationnelle~; le retour de l'État comme financeur de premier rang (CHIPS Act) reproduit un pattern documenté depuis~1838.

\textcolor{teal}{Ce fait stylisé ne mentionne pas l'énergie comme objet. L'énergie intervient comme composante du coût de fonctionnement permanent qui distingue les infrastructures actives (data centers, qui consomment en continu) des infrastructures passives (fibre optique posée, câble sous-marin abandonné, carte perforée dans sa boîte). Cette distinction est constitutive du type de capital (D3)~; ses conséquences pour le financement sont traitées au chapitre~3. Ce que D8 établit en propre, c'est que la structure de financement est un déterminant du déploiement au même titre que la structure technique (D1--D3) et la structure de marché (D2, D6, D7).}


% =============================================================================
